\chapter{Cengage DPP (3.1 to 3.5): Complete Question Bank}
\label{chap:cng_dpp_all_questions}

\begin{problembox}
\textbf{Problem 1} \hfill \textbf{[Cengage DPP Q9]}
\label{q:ext-cng_dpp_all_questions-1}

The equilibrium constant of acetic acid in an aqueous
your answer.
solution of concentrationcis given by
3.Why on dilution the Am of CHCOOHincreases drastically,
CA
while that of CH,COONa increases gradually?
CA2
=X (e)
(b)K=
V-$/infty$V
($2$V-V)V
4.A 100 watt, 110 volt incandescent lamp is connected
inseries with an electrolytic cell containing cadmium
CA2
CAc
sulphate solution.Whatmass.of cadmium willbe deposited
(c)K=
(d)K=
$2$v+V
A$/circ$(A$/infty$-A)
by theflowingfor10 hours?
\end{problembox}

\begin{problembox}
\textbf{Problem 2} \hfill \textbf{[Cengage DPP Q5]}
\label{q:ext-cng_dpp_all_questions-2}

The resistance of a solutionA'is 50 $/Omega$ and that of
\end{problembox}

\begin{problembox}
\textbf{Problem 3} \hfill \textbf{[Cengage DPP Q10]}
\label{q:ext-cng_dpp_all_questions-3}

The molar conductivity of 0.05 M of solution of an
solution'B'is100 $/Omega$,both solutions being takenin the
electrolyte is 200 Q- cm$2$mol-'. The resistance offered
same conductivity cell. If equal volumes of solutions A
by a conductivity cell with cell constant (1/3) cm-/ would
andBaremixed,whatwillbetheresistance of themixture
be about
using the same cell?(Assume that there is no increase in
(a) 11.11 Q
(b)22.22Q
the degree of dissociation of A and B on mixing).
(c)33.33$/Omega$
(d) 44.44 Q
6.The equivalent conductivity of acetic acid at infinite
\end{problembox}

\begin{problembox}
\textbf{Problem 4} \hfill \textbf{[Cengage DPP Q11]}
\label{q:ext-cng_dpp_all_questions-4}

If 2m (Kt) = 73.5 Q-l cm$2$ mol-1. /textasciitilde()m (At) = 189 Q-1
dilution is 387 S cm$2$ eq-'.At the same temperature,
0.001 M solution of acetic acid,it is 55 S cm$2$eq-'. What
cm$2$mol-l and (SO$2$) = 160 Q-cm$2$mol-, the molar
is the degree of dissociation of 0.1N acetic acid?Assume
conductivity of potash alum will be about
1-$/alpha$=1.
(a) 1165 Q-l cm$2$ mol-l
7.Calculate the dissociation constant of water at 25$/circ$Cfrom
(b)422.5 Q-1 cm$2$ mol-1
the following data:
(c)582.5 Q-l cm$2$mol-l
Conductivity of H0 = 5.8 x 10-8 S cm-1, /textasciicircum()g+ =350.0
(d)145.6 Q-1 cm$2$mol-1
and 2oH- =198.0S cm2eq.-!
--- PAGE 2 ---
3.2Physical Chemistry
12.The molar conductivity of a weak electrolyte HA at infinite
dilution can bedeterminedby
(c)
(p)
(a)extrapolating $/Delta$m versus c plot to c $/to$ 0.
(b)extrapolating $/Delta$m versus $/sqrt()$c plot to c $/to$ o.
(c)using Kohlrausch law with the data on
MultipleCorrectAnswersType
A(HCI), A (NaA) and A(NaCI)
(d)using Kohlraush law with the data on.
18.Which one is a correct statement?
(a)Conductance is an additive property
A(HA) and A(NaA).
(b) Conductance order: Litaq)<Naa)<Kaq)< Rb
13.When a certain conductivity cell was filled with 0.01 M
(c)Molar conductivity of NH in liquid NH is highe
solution of KCl, it had a resistance of 160Q at 25$/circ$C,and
when filled with 0.005 M NaOH, it had a resistance of
than in H2O
190 Q.If specific resistance of KCl solution is 700 Q-cm,
(d)At infinite dilution all the interionic attraction
specific conductance (Q-1 cm-1) solution is
vanishes
(a)0.00120
(b)0.00170
19.An ion travels 0.4 cm in1 minute under the influence
(c)0.00180
(d)0.00190
2 volt potential difference when electrodes are 3 cm apa
14.Variation of resistance with increase incell constant gives
Which one are correct?
graph of the type
(a) Speed of ion = 6.67 $/times$ 10-3 cm s-1
(b)Ionicmobility =1x10-6m$2$volt-1 sec-1
R
R
(c).Ionic mobility = 1.0 $/times$10-3 cm$2$ volt1 sec-1
(d)Ionicmobility = 4.44 $/times$10-3 cm$2$volr-  sec-1
(a)
(b)
20.Which ion(s) has/have exceptionallyhigher Avalues?
(a) H+
(b)K+
[|Cellconstant
(c)OH-
(d) NH2
21.For which,the formula $/alpha$=
 . 
(a) CH,COOH
(b) HCl
R
(c) HCiO4
(d) CH,NH,
(c)R
(d)
22.The factor(s) which influence(s) the conductance
solutions is.
(a)Solute-solute interaction
(b)Solute-solvent interaction
(c)Temperature
15.The molar conductivity of LiNO at infinite dilution is
(d)Solvent-solventinteraction
105S cm$2$mol-1.If 0.02 M solution shows conductance
23.Which of the following graph incorrectly representst
of 0.158 $/Omega(-1)$ using cell of cell constant 1 cm-'. The
variation of molar conductance ($/Lambda$m)with dilution fol
concentration of Lit ion in g litreis:
strong electrolyte?
(a) 0.01106
(b) 0.1106
(c) 11.06
(d)0.2212
\end{problembox}

\begin{problembox}
\textbf{Problem 5} \hfill \textbf{[Cengage DPP Q16]}
\label{q:ext-cng_dpp_all_questions-5}

The absolute velocity of Ag+ and NO are 5.7 x
$/sqrt()$(e)
(b) $/Lambda$m
10-4 cm sec-1 and 6.9 $/times$ 10-4 cm sec-1 respectively.
Assumingcompletedissociation ofAgNO,calculate the
Dilution-
Dilution-
molar conductivity of AgNO at infinite dilution.
(a)121.59 S cm$2$ mol-1
(b)141.59 S cm$2$mol-l
(c)131.59 S cm$2$mol-'(d) none of these
(c) $/Delta$m
(d)Am
and the conductivity of the solution was measured. The
plots of conductanceA vs volume of AgNO(aq)is:
Dilution---
Dilution---$/to$
a)
v (q)
\end{problembox}

\begin{problembox}
\textbf{Problem 6} \hfill \textbf{[Cengage DPP Q24]}
\label{q:ext-cng_dpp_all_questions-6}

--- PAGE 3 ---
Electrochemistry 3.3
Comprehension Type
Aboveplotrepresents thevariationofmolarconductance
against $/sqrt()$C(where C =molar concentration of the
The molarconductance of NaCl varies with theconcentration as
electrolyte). Select the incorrect options among the
shown in the following table:
following:
Molar conductance in
(a).Both I and II are for strong electrolyte
Molar concentration
(b)Both I andIlare forweak electrolyte
$/Omega(-1)$cm$2$mol-1
of NaCI
()Iisforstrong electrolyte andIIforweak clectrolyte
107
4 $/times$ 10-4
(d)Iis for weakelectrolyte and II for strong electrolyte
97
9$/times$10-4
MatchingColumnType
87
16$/times$10-4
\end{problembox}

\begin{problembox}
\textbf{Problem 7} \hfill \textbf{[Cengage DPP Q25]}
\label{q:ext-cng_dpp_all_questions-7}

Match the physical term in Column I with the unit in
Whenacertain conductivity cell(C) was filled with 25$/times$10-4(M)
Column Il.
NaCl solution, the resistance of the cell wasfound tobe1000
ColumnII
$/Omega$.At infinite dilution ion conductances of Cl- and SO2 are
ColumnI
80 $/Omega(-1)$ cm2mol-1 and 160 $/Omega$-l cm2mol-l respectively.
(a)
Conductance
(p)
siemen m$2$mol-1
28.Whatis themolarconductance of NaClat infinite dilution?
(b)
Resistivity
(q)
m-1
(a)147 $/Omega(-1)$ cm$2$mol-1
(c)
Conductivity
(1)
volt perampere
(b)107 $/Omega(-1)$ cm$2$ mol-1
(d)
Molar conductivity
(s)
siemen
(c)127 $/Omega(-1)$ cm$2$mol-1
Cell constant
(/\#)
siemenm-1
(d)157$/Omega$-cm$2$mol-1
(e)
29.What is the cell constant of the conductivity cell(C)?
(f)
Resistance
(u)
$/Omega$metre
(a)0.385 cm-1
(b)3.85 cm-1
(c)38.5 cm-l
(d) 0.1925 cm-1
IntegerType
30.If the cell(C) is filledwith 5$/times$10-3(N)NaSO4,the
observed resistancewas 400 $/Omega$.What is the molar
26.The molar conductivity of acetic acid at infinite dilution
conductance of NaSO?
is 390.7and for 0.01M acetic acid is 3.907 S cm$2$ mol-
(a)19.25 $/Omega$-l cm$2$mol-1
What is the pH solution?
(b)96.25 $/Omega$-l cm$2$ mol-1
27.The specific conductivity of N/50 solution of KClat298K
(c)385$/Omega(-1)$cm$2$mol-1
is 0.002501S cm-1.If the resistance of the same solution
(d)192.4 $/Omega(-1)$ cm$2$mol-1
placed in the cellis 2000 $/Omega$s, what is the cell constant?
AnswersKey
SubjectiveType
MatchingColumn Type
\end{problembox}

\begin{problembox}
\textbf{Problem 8} \hfill \textbf{[Cengage DPP Q4]}
\label{q:ext-cng_dpp_all_questions-8}

19.06 g 5. 66.67 $/Omega$
(1$/leftarrow$J)(b$/leftarrow$)(d$/leftarrow$p)(1$/leftarrow$o)(n$/leftarrow$q)s$/leftarrow$)s
6.1.5/%
7.2$/times$10-16 mol/litre
IntegerType
SingleCorrectAnswerType
\end{problembox}

\begin{problembox}
\textbf{Problem 9} \hfill \textbf{[Cengage DPP Q12]}
\label{q:ext-cng_dpp_all_questions-9}

(c)
Comprehension Type
\end{problembox}

\begin{problembox}
\textbf{Problem 10} \hfill \textbf{[Cengage DPP Q30]}
\label{q:ext-cng_dpp_all_questions-10}

(d)
MultipleCorrectAnswersType
\end{problembox}

\begin{problembox}
\textbf{Problem 11} \hfill \textbf{[Cengage DPP Q24]}
\label{q:ext-cng_dpp_all_questions-11}

(a, b, d)
--- PAGE 4 ---
DPP
Electrolysis and Faraday's Laws
9.Which of the following gives the charge of an electron?
SubjectiveType
(a) 1 F/1 NA
(b) 1F
(c)1F$/times$1N
(d)1N/1F
\end{problembox}

\begin{problembox}
\textbf{Problem 12} \hfill \textbf{[Cengage DPP Q1]}
\label{q:ext-cng_dpp_all_questions-12}

An aqueous solution of copper sulphate is electrolysed
using platinum electrodes inone case and copper electrodes
10.In the electrolysis ofcopper(H)chloridesolution,the mas
of cathode is increased by3.15g.Atcopper anode,we will
in another case.Will the products ofelectrolysis besame
or different?Give reason.
have
\end{problembox}

\begin{problembox}
\textbf{Problem 13} \hfill \textbf{[Cengage DPP Q2]}
\label{q:ext-cng_dpp_all_questions-13}

State the products of electrolysis obtained on the cathode
(a)the liberationof1120mLof Clat STP
and the anode in the following cases:
(b) the liberation of 560 mL of O2 at STP
(i)A dilute solution of H,SO4 with platinumelectrodes.
(c)the loss of 3.15 g of Cu
(i)An aqueous solution of AgNO with silver electrodes.
(d) the gain of 3.15 g of Cu
3.A current of 2 ampere is passed for one hour between nickel
11.Onegequivalent of Nametal is formed from electrolysis
electrodes in0.5 L of 2 M Ni(NO) solution.What will
of fusedNaCl.No.of moles of Al from the fused NaAIF
be themolarityof the solution at the end ofelectrolysis?
with the same current passed is:
4.(i) On electrolysis of an aqueous solution NaCl, why H2
(a1
(b)3
(c)1/3
(d) 2
and not Na is liberated at the cathode?
(i) One faraday of electricity deposits one mole of
MultipleCorrectAnswersType
Na from the molten salt but
molofAlfroman
12.Which of the following statements are correct?
aluminium salt.Why?
3
(a)Sodium canbeobtainedbytheelectrolysisof aqueous
solutionof NaClusingPtelectrodes.
Single CorrectAnswerType
(b)The current carryingionsin anelectrolytic cell arn
not necessarily discharged at the electrodes.
5.In electrolysis of an aqueous solution of sodium sulphate,
(c)Themass of a substance deposited on the cathode
2.4 L of oxygen at STPwas liberated at anode.Thevolume
anode during electrolysis is given asm=QM/FIv
of hydrogen at STPliberated at cathode would be:
(d)In the purification of a metal by electrolytic method
(a) 1.2 L
(b)2.4 L
the impure metalis made cathode.
(c)2.6L
(d)4.8L
13.Which of the following statements are not correct?
6.What mass of silvercould be plated out on a serving trayby
.(a)Thereaction involvingreductionofwateri
electrolysisof a solution containing silverin+1 oxidation
2HO+2e$/to$H+2OH
statefor a period of 8hours at a current of 8.46 amperes?
(b) The species having a least reduction potentiali
What is the areaof trayif the thickness of the silverplating
preferentiallyreduced at the cathode.
is 0.00254 cm? Density of silver is 10.5 g/cm3.
(c)Water is reduced in preference to the cations havin
(a)11200 cm$2$
(b)10204.7cm$2$
morepositivepotential than-0.828V.
(c)10300cm$2$
(d)9237.4 cm2
(d)In the electrolysis of a solution containing SO$2$ anion
\end{problembox}

\begin{problembox}
\textbf{Problem 14} \hfill \textbf{[Cengage DPP Q7]}
\label{q:ext-cng_dpp_all_questions-14}

A certain electricity deposited 0.54 g of Ag from AgNO3
only, the anodic reaction is 2HO+ 2e-$/to$H+20H
solution.What volumeof hydrogenwill the same quantity
\end{problembox}

\begin{problembox}
\textbf{Problem 15} \hfill \textbf{[Cengage DPP Q14]}
\label{q:ext-cng_dpp_all_questions-15}

Which of the following statements are not correct?
of electricity liberate at 27$/circ$C and 728mm Hg pressure?
(a)The electrolysis of a solution containing Cl aniom
(a) 62.48 mL
(b)60.00 mL
produces Cl, at anode but O is also liberated as th
(c)64.28 mL
(d) 66.5-mL
solutionisdiluted.
8.The time taken to decompose electrolytically18gwater
(b)The quantitative relationship between electricity an
by 2A current is about
chemical change were established byFaraday.
(a) 24.8 h
(b)26.8h
(c) When an aqueous solution of sodium fluoride
(c) 28.8 h
(d) 30.8 h
electrolysed, the gas liberated at the anode is F
(d) In anelectrolyticrefining of an impure metal, thpu
metal is obtained at anode.
--- PAGE 5 ---
Electrochemistry 3.5
19.A current strength of 96.5 A is passed for 10 sec through
\end{problembox}

\begin{problembox}
\textbf{Problem 16} \hfill \textbf{[Cengage DPP Q15]}
\label{q:ext-cng_dpp_all_questions-16}

Three molesofelectrons are passed through three solutions
1 L of solution of 0.1M aqueousCuSO,calculate the pH
in succession containing AgNO, CuSO4 and AuCl,
of the solution.
respectively.Themolarratioofamounts ofcationsreduced
at cathode will be
ComprehensionType
1.11
(a) 1:2:3
(b)
123
Electrolysisisthe process inwhichelectricalenergyisconverted
(c) 3:2:1
(d) 6:3:2
to chemical energy.In electrolytic cell,oxidation takes place at
16.On electrolysis, in which of the following, O2would be
anode and reduction at cathode.Electrode process dependson the
liberated at the anode?
electrode takenfor electrolysis.Amount ofsubstanceliberated
(a) dilute H,SO4with Pt electrodes
at an electrode is directlyproportional to the amount of charge
(b)aqueous AgNO, solution with Pt electrodes
passed through it.The mass of substance liberated at electrode
(c)dilute H,SO with Cu electrodes
is calculated using thefollowingrelation: 
(d)aqueous NaOH with aFe cathode and a Pt anode
ItE
m=
96500
Matching ColumnType
Here,Erepresents theequivalent mass and96500 Cis called
theFaraday constant.Faraday(96500 C)is the charge of 1 mol
17.Match the Column I with Column II:
electron,i.e,6.023$/times$102electrons;it is used toliberate one
Column I
ColumnII
gram equivalent of the substance.
(Electrolysis)
(pH and products at
20.The platinum electrodes were immersed in a solution of
anode and cathode)
cupricsulphate(CuSO) and electric current is passed
Electrolysis of 100 L
(p)
pH =12.3
(a)
throughthesolution.Aftersometime,itwasobservedthat
aqueous solution
Anode =Ethane(g)
thecolour of coppersulphatedisappearedwithevolution
+CO(g)
CHCOOK by passing
of a gas at the electrode.The colourless solution contains:
2Fof electricity
Cathode =H(g)
(a)platinum sulphate
(b)copper nitrate
(c)copper sulphate
(d) sulphuric acid
(b)
Electrolysis of 10 L
(b)
pH=13.0
aqueous solution
Anode = H(g)
21.The passage of current librates H,at cathode and Cl,at
HCOOK by passing
+CO(g)
anode.The solution is:
1 Fofelectricity
Cathode=H(g)
(a) copper chloride in water
(b)NaCl in water
(c)
Electrolysis of 10 L
(1)
pH=7.0
(c)mercuric chloride in water
aqueous solution
Anode=O(g)
Cathode = H(g)
(d)AuCl in water
K,SO4 by passing 1 F
\end{problembox}

\begin{problembox}
\textbf{Problem 17} \hfill \textbf{[Cengage DPP Q22]}
\label{q:ext-cng_dpp_all_questions-17}

On electrolysis of dilute sulphuric acid using platinum
ofelectricity
electrodes,the product obtained at the anode will be:
(d)
Electrolysis of 10 L
(s)
pH=1.0
(a)hydrogen
(b) oxygen
aqueous solution
Anode= O(g)
(c)hydrogen sulphite
(d)sulphur oxide
CuF2 by passing 1 F
Cathode = Cu(g)
23.How many faradays are required to reduce 1 mol BrO; to
ofelectricity
Br-?
(a)3
(b) 5
(c) 6
(d) 4
24.Calculate the volume of gas liberated at the anode at S.T.P.
IntegerType
during the electrolysis of a CuSO solution by a current of
1 A passed for 16 min and 5 sec.
\end{problembox}

\begin{problembox}
\textbf{Problem 18} \hfill \textbf{[Cengage DPP Q18]}
\label{q:ext-cng_dpp_all_questions-18}

How many grams of Cl(g) can be produced by the
(a)224 mL
(b)56 mL
electrolysis of molten NaCl with a current of 5.5Afor 25
(c)112mL
(d) 448 mL
min?[Atomic weight of Cl = 35.5 amu]
AnswersKey
Matching Column Type
Single Correct Answer Type
\end{problembox}

\begin{problembox}
\textbf{Problem 19} \hfill \textbf{[Cengage DPP Q17]}
\label{q:ext-cng_dpp_all_questions-19}

(a $/to$ p),(b $/to$q),(c $/to$r),(d $/to$s)
5.(d)
10.(c)
\end{problembox}

\begin{problembox}
\textbf{Problem 20} \hfill \textbf{[Cengage DPP Q11]}
\label{q:ext-cng_dpp_all_questions-20}

(c)
IntegerType
Mutiple Correct Answers Type
\end{problembox}

\begin{problembox}
\textbf{Problem 21} \hfill \textbf{[Cengage DPP Q14]}
\label{q:ext-cng_dpp_all_questions-21}

(c, d)
Comprehension Type
\end{problembox}

\begin{problembox}
\textbf{Problem 22} \hfill \textbf{[Cengage DPP Q24]}
\label{q:ext-cng_dpp_all_questions-22}

(b)
--- PAGE 6 ---
3.6Physical Chemistry
DPP 3.3
Fuel Cell and Batteries
(a)Pb2++2e$/to$Pb
SubjectiveType
(b)Pb$/to$Pb2++ 2e-
(c)Pb2++SO$2$-$/to$PbSO4
1.(i)What advantage do the fuel cells have over primary
(d)PbSO4+2H2O$/to$PbO2+4H++SO2-+2e
and secondary batteries?
(i)Write the cellreaction ofalead storagebattery whenit
10.Hydrogen-oxygen fuel cells are used in spacecraft to
is discharged.How does the density of the electrolyte
supply:
change when the battery is discharged?
(a)power for heat and light
\end{problembox}

\begin{problembox}
\textbf{Problem 23} \hfill \textbf{[Cengage DPP Q2]}
\label{q:ext-cng_dpp_all_questions-23}

(i) What is the role of ZnCl, in a dry cell?
(b)power for pressure
(c)oxygen
its life?
(d)none of these
3.Which type of cells are rechargeable?
11.In a H-Ocell,combustion of hydrogen occurs due to:
4.During the discharge of a lead storage battery, density of
(a)remove adsorbedoxygenfromelectrodesurface
H,SO fell from 1.294 g/mL to 1.139 g/mL.Sulphuric acid
(b)create potential differencebetweenthe twoelectrodes
of density1.294 g/mL is 39/%H,SO4by weight and that of
(c)produce highpuritywater
density 1.139 g/mL is 20/% HSO4by weight.The battery
(d)generateheat
holes 3.5 litres of the acid andvolume remains practically
constant duringdischarge.Calculate ampere-hour of
MultipleCorrectAnswersType
which the batterymust have been used.The charging and
discharging reactions are:
\end{problembox}

\begin{problembox}
\textbf{Problem 24} \hfill \textbf{[Cengage DPP Q12]}
\label{q:ext-cng_dpp_all_questions-24}

Which is/are correct about fuel cells?
Pb+SO$2$-$/to$PbSO+2e
(charging)
(a) Cells continuously run as long as fuels are supplied
(b) These are more efficient and free from pollution
PbO2+4H*+SO$2$+2e $/to$PbSO+2H2O(discharging)
(c)These are used toprovidepower and drinkingwater
5.In a fuel cell H,and Oreact to produce electricity.In the
toastronautsinspaceprogramme
processHgasis oxidised at the anode and Oisreduced at
(d)None of these
the cathode.If 67.2 litres of Hat NTPreacts in15 minute,
\end{problembox}

\begin{problembox}
\textbf{Problem 25} \hfill \textbf{[Cengage DPP Q13]}
\label{q:ext-cng_dpp_all_questions-25}

During the discharge of lead storage battery:
what is the average current produced?If the entire current
(a)PbOis reduced
is used for electro-deposition of Cu from Cu$2$+,how many
(b)Pb is oxidised
gof Cu are deposited?
(c)specific gravity of HSOgives extent of discharging
(d)equivalent mass of HSOis 98
\end{problembox}

\begin{problembox}
\textbf{Problem 26} \hfill \textbf{[Cengage DPP Q14]}
\label{q:ext-cng_dpp_all_questions-26}

Which of the following statements are correct about Ni-Cd
SingleCorrectAnswerType
cell
6.When a lead storage battery is discharged,then:
(a). It consists of a cadmium electrode (as anode) and a
(a)SOis evolved
metal grid containing nickel (IV) oxide (as cathode)
(b) lead is formed
immersed in KOH solution.
(c) lead sulphate is consumed
(b)At cathode:NiO(s)+2HO+2e=Ni(OH)(s)+
(d)sulphuric acid is consumed
2OH-(aq)
7.A depolarizer used in dry cell batteries is:
(c)At anode:Cd(s)+ 2OH-(aq)=Cd(OH)(s)+ 2e
(a)ammonium chloride (b)manganese dioxide
(d)The potential of each Ni-Cd cellis approximately 2.6
(c)potassium hydroxide (d) sodium phosphate
V.
8.An example of a simplefuel cell is:
\end{problembox}

\begin{problembox}
\textbf{Problem 27} \hfill \textbf{[Cengage DPP Q15]}
\label{q:ext-cng_dpp_all_questions-27}

During the charging oflead storage battery, which reactins
(a)lead storage battery
(b)H2-O2cell
don't occur at anode are:
(c)Daniell cell
(d)Leclanche cell
(a)Pb$2$++SO$2$-$/to$PbSO4
\end{problembox}

\begin{problembox}
\textbf{Problem 28} \hfill \textbf{[Cengage DPP Q9]}
\label{q:ext-cng_dpp_all_questions-28}

Which of the following reactions occurs at cathode during
:charging of storage battery?
(b)PbSO4+H2O$/to$PbO2+SO2-+2H+
--- PAGE 7 ---
Electrochemistry 3.7
ComprehensionType
(c)Pb$/to$Pb$2$++2e
(d)Pb$2$++2e$/to$Pb
A lead storage battery consists of a lead anode and a grid of
\end{problembox}

\begin{problembox}
\textbf{Problem 29} \hfill \textbf{[Cengage DPP Q16]}
\label{q:ext-cng_dpp_all_questions-29}

Fuel celinvolves following reaction(s):
(at cathode)
lcad packed withlead dioxide as the cathode.The electrolyte
(a) O(g) + 2HO(1) + 4e $/to$ 4OH- (aq.)
(at anode)
taken is 39/% HSO4 by mass having a density of I.294 g mL-1.
(b)O(g)+ 2HO(1) + 4e $/to$ 4OH-(aq.)
The batteryholds 35Lof the acid.During the discharge of the
(at anode)
(c)2H(g)+ 4OH-(aq)$/to$4HO(1)+ 4e
battery,the density H2SO4 falls from 1.294 g mL- to 1.139 g
(at cathode)
(d)2H(g)+4OH-(aq)$/to$4HO(1)+ 4e
mL-I which is 20/% HSO4 by mass.
\end{problembox}

\begin{problembox}
\textbf{Problem 30} \hfill \textbf{[Cengage DPP Q19]}
\label{q:ext-cng_dpp_all_questions-30}

The reaction occurring at the anode during charging is
MatchingColumnType
(a)Pb2++2e$/to$Pb
(b)Pb$2$++SO2$/to$PbSO4
17.Match the following:
(c)Pb$/to$Pb2++2e
Column II
(d)PbSO4+2H2O$/to$2PbO2+4H++SO$2$-+2e
ColumnI
(p)
Cellreaction 2H+O2$/to$2HO
20.Moles of sulphuric acid lost during discharge is
(a)
Leclanche
(a) 9.88
(b) 8.88
cell
(c) 7.88
(d) 6.88
(b)
Ni-Cd cell
(b)
Does not involve any ion in
21.Molarity of the solution after the dischargeis
solution and is used in hearing
(a) 8.136
(b) 4.068
aids.
(c)2.32
(d) 1.16
(c)
Fuel cell
(r)
Rechargeable
22.The amount of charge in colulombs used upby the battery
Mercury cell
(s)
Reaction at anode,
is nearly
(d)
(a)954180
(b)477090
++uzuz
(c)95418
(d) 47709
23.The number of ampere-hour for which the battery must
IntegerType
havebeen used is
(a)2650.5
(b) 265.05
\end{problembox}

\begin{problembox}
\textbf{Problem 31} \hfill \textbf{[Cengage DPP Q18]}
\label{q:ext-cng_dpp_all_questions-31}

The emf of lead storage battery is .... V.
(c)26.505
(d)2.6505
AnswersKey
MatchingColumnType
SubjectiveType
\end{problembox}

\begin{problembox}
\textbf{Problem 32} \hfill \textbf{[Cengage DPP Q17]}
\label{q:ext-cng_dpp_all_questions-32}

(a $/to$s),(b $/to$r),(c $/to$p),(d $/to$q)
4.265.04 5. 643.3 amp,190.5 g
IntegerType
Single CorrectAnswerType
9.(a)
\end{problembox}

\begin{problembox}
\textbf{Problem 33} \hfill \textbf{[Cengage DPP Q11]}
\label{q:ext-cng_dpp_all_questions-33}

(b)
Comprehension Type
\end{problembox}

\begin{problembox}
\textbf{Problem 34} \hfill \textbf{[Cengage DPP Q23]}
\label{q:ext-cng_dpp_all_questions-34}

(b)
MultipleCorrectAnswersType
\end{problembox}

\begin{problembox}
\textbf{Problem 35} \hfill \textbf{[Cengage DPP Q16]}
\label{q:ext-cng_dpp_all_questions-35}

(a, c)
--- PAGE 8 ---
DPP
Construction and Working of a Cell, Electrochemical Series and
Its Applications
Which is the strongest reducing agent?
SubjectiveType
(a) Zn(s)
(b) Cr(s)
(c)H(g)
(d)Fe2+(aq)
1.Predict whether zinc and silver react with 1M sulphuric
6.The following facts are available:
acid togive outhydrogengasornot.Giventhat the standard
2X-+Y2$/to$2Y-+X2;
reduction potentials of zinc and silver are - 0.76 volt and
2W- + Y2 $/to$No reaction;
0.80 volt respectively.
2Z- +X22x-+Z2
2.Iodine (I) and bromine (Br)are added to a solution
containing iodide (I-)and bromide(Br)ions.What
Whichof the followingstatementsis correct?
reaction would occur if the concentration of eachspecies
(a)
FO
is1 M?The electrode potentials for the reactions are:
W-/W
E=0.54V,E
Br/Br-=1.08V
(b)
1/
W/W2
3.(i)Which of the following oxides is reduced by hydrogen?
(c)
W/W
MgO,CuO and NaO
(ii)Which of thefollowing oxides.will decompose on
W/W2
Z/Z2
heating?
AgasX at1atm isbubbled througha solution containing
ZnO,CuO,MgO and Ag2O
a mixtureof1MYand1MZat25$/circ$C.If thereductior
(ii)The value of E for electrode reactions,
potential of Z>Y>X,then
(a)Y will oxidise X and not Z
+ 2e are 0.444,-0.337 and 0.763 volts respectively.
(b)Ywill oxidise Zand not X
Statewhichof thesemetalscanreplacetheothertwo
(c)Y will oxidise both X and Z
from the solution of their salts?
(d)Y will reduce both X and Z
4.DeterminerangeofE$/circ$valuesforthisreaction
8.Alargequantity ofcobalt metal was droppedintoa solutio
x2+ + 2e $/to$ X(s)for given conditions:
containing Ag+,Fe3+and Cu$2$+,all at unit concentration.
E$/circ$(Ag+/Ag)=0.799V,E$/circ$(Fe3+, Fe2+/Pt)=0.77V,E(Cu$2$
(a)If themetalXdissolvesinHNObut notinHClit can
Cu)=0.34V,E$/circ$(Fe3+/Fe)=-0.04V,E$/circ$(Cu$2$+,Cu+/Pt)
displace Ag+ ion but not Cu2+ion.
0.153 V and E$/circ$(Co2+/Co)=0.28V,then the solution woull
(b)If the metal X is HCl acid producing H(g)but does
contain
not displaceeitherZn+orFe2+
Given:
(a)Ag,Fe and Cu
(b)Ag,Fe and Cu+
(c)Fe2+, Cu and Co2+
(d)Fe$2$+,Cu,Ag and Co+
=0.8V, 
0.44V, E
= 0.34V, 
C2+/C
9.Given :E(Sn$2$+,Sn++/Pt)= 0.15 V,E$/circ$(Hg$2$. Hg-+/Pt)
-0.76V
0.92 V and E$/circ$(Pb2+, H+/PbO2)= 1.45 V. Based on this dan
=0.96V
NONO
Zn$2$+/Zn
which of thefollowing statements is correct?
(a)Sn4+ is a stronger oxidizing agent than Pb++
Single CorrectAnswerType
(b) Sn2+ is a stronger reducing agent than Hg?*
(c)Hg2+ is a stronger oxidizing agent than Pb++
5.The standard potential at 298 K for the following half
(d)Pb2+ is a stronger reducing agent than Sn+
reactions are given against each:
10.Given below are some standard reduction potentials:
Zn$2$(aq)+2eZn(s)
0.762 V
Sn$2$++2e=-0.14V
Hg+2e-=2HgE=0.79
2H+(aq)+ 2eCr(s)
0.000V
Which of the following statements is correct?
Cr3+(aq)+3eCr(s)
0.740 V
(a)Hg can reduce Sn$2$+ to Sn
Fe3+(aq) + 2eFe2+(aq)
0.770V
(b) Sn2+ can oxidize Hg to Hg2+
--- PAGE 9 ---
Electrochemistry 3.9
\end{problembox}

\begin{problembox}
\textbf{Problem 36} \hfill \textbf{[Cengage DPP Q19]}
\label{q:ext-cng_dpp_all_questions-36}

In which of the following,salt bridge is not needed?
(c)Hg2+ can oxidize Sn to Sn2+
(a) PbIPbSO4(s)I H2SO41 PbO2(s)I Pb
(d) Between Hg2 and Sn2+, Sn2+ is stronger oxidising
(b) Cd I CdO(s)I KOH(aq)I NiO2(s)INi
agent than Hg2+.
(c)Fe(s)IFeO(s)I KOH(aq),NiO(s)INi
11.Given that E values ofAg+/Ag,K+/K,Mg+/Mg and Cr3+/
(d)ZnIZnSOICuSOICu
Cr are 0.80V,-2.90V,-2.37V,and-0.74V,respectively.
Whichof the following orders.regarding thereducing
MatchingColumnType
power of metals is correct?
(a)Ag>Cr> Mg>K
(b)Ag<Cr<Mg<K
20.Match the items of ColumnI and Column II on the basis
(c)Ag>Cr>K>Mg
(d) Cr>Ag>Mg >K
of data given below:
\end{problembox}

\begin{problembox}
\textbf{Problem 37} \hfill \textbf{[Cengage DPP Q12]}
\label{q:ext-cng_dpp_all_questions-37}

Indicator clectrode is:
EO
=2.87V, 
E
=-3.5V
(b) Calomel electrode
F2/F
Li/Li
(a)SHE
(c)Ag/AgClelectrode
(d)Quinhydrone electrode
EO
=1.4 V,E
=1.09 V
Au3+/Au
Br/Br
13.When an aqueous solution of lithium chloride is
electrolysedusing graphite electrodes:
ColumnI
ColumnII
(a)pH of the resulting solution increases
(b)pH of the resulting solution decreases
(a)
F2
(P)
Metal is the strongest reducing agent
(c)as the current flows,pH of the solution around the
(b)
Li
(q)
Metalionwhichis the weakest
cathode increases
oxidisingagent
(d)none of the above
14.The calomel electrode is reversible with.respect to:
(c)
Au3+
(r)
Non-metal which is the best oxidising
(a) Hg2+
(b)H+
(c)Hg$2$+(d) CI-
agent
(d)
Br
(s)
Unreactive metal
MultipleCorrectAnswersType
(e)
Au
()
Anion that can be oxidised by Au3+
15.Which of the following is/are true for standard hydrogen
(f)
Lit
(u)
Anion which is the weakest reducing
electrode?
agent
(a)Temperature is maintained at 25$/circ$$/circ/text(C)$
(g)F-
(v)
Metal ion which is an oxidising agent
(b)Hydrogen ion concentration is1 M
(c)Pressure of hydrogengasismaintained at 1atmosphere
(d)It contains platinum metal which is used to adsorb
IntegerType
hydrogen gas, 
\end{problembox}

\begin{problembox}
\textbf{Problem 38} \hfill \textbf{[Cengage DPP Q16]}
\label{q:ext-cng_dpp_all_questions-38}

The electrode potential of a glass electrode does notdepend
\end{problembox}

\begin{problembox}
\textbf{Problem 39} \hfill \textbf{[Cengage DPP Q21]}
\label{q:ext-cng_dpp_all_questions-39}

If E$/circ$,E2 and E are standard oxidation potentials for
:uodn
(a)concentration of chlorideions
The value of n is .....
n
(b) concentration of hydrogen ions
\end{problembox}

\begin{problembox}
\textbf{Problem 40} \hfill \textbf{[Cengage DPP Q22]}
\label{q:ext-cng_dpp_all_questions-40}

The no. of cells which may be constructed with different
(c)concentration of KCl solution
(d) all of these
ECen values for the reaction: Fe + 2Fe3+ $/to$ 3Fe2+.
17.A KI solution containing starch,turns blue on the addition
of chlorine.Which of the following statements explain(s)
Comprehension Type
this?
(a)The standard reduction potential of Cl,is more than
The standard reduction potentials of some half reactions are
that of L2
given as below:
(b)The standard oxidation potential of Cl, is more than
Cu$2$+e$/to$Cu*;E$2$=0.15V
that of I
Cu++e$/to$Cu;E2=0.5V
(c)The product formed when Cl, combines with starch
is blue in colour
Ag++e$/to$Ag;E$2$=0.799V
(d))The produet formed when I combines with starch is
The decrease in free energy (-$/Delta$G) is responsible .for the
bluein colour
production of equivalent amount of electrical work. Thus,
\end{problembox}

\begin{problembox}
\textbf{Problem 41} \hfill \textbf{[Cengage DPP Q18]}
\label{q:ext-cng_dpp_all_questions-41}

Which of the following statements is/are correct?
(a)Ametalinits highest oxidationstateacts as anoxidant
-$/Delta$G=nFxE
\end{problembox}

\begin{problembox}
\textbf{Problem 42} \hfill \textbf{[Cengage DPP Q23]}
\label{q:ext-cng_dpp_all_questions-42}

The ECen for CuI Cu2+ in volt is:
(b)In the reaction,F2+O2
$/to$OF,oxygen is an
2
(a) 0.35
(b)-0.35(c)+0.325 (d)-0.325
oxidant
(c)The oxidant number of Ni in Ni(CO)g is zero
\end{problembox}

\begin{problembox}
\textbf{Problem 43} \hfill \textbf{[Cengage DPP Q24]}
\label{q:ext-cng_dpp_all_questions-43}

The Een for CuI Cu2+IlAg+IAg in volt is:
(d)Copper metal can be oxidized by Zn2+ ions
(a) +0.449
(b)+0.474
(c)-0.449
(d)-0.474
--- PAGE 10 ---
3.10Physical Chemistry
(a) 1,3
(b) 2,3
\end{problembox}

\begin{problembox}
\textbf{Problem 44} \hfill \textbf{[Cengage DPP Q25]}
\label{q:ext-cng_dpp_all_questions-44}

The Een for Cu|Cu2IAg*IAg is:
(c) 1,4
(d) 2, 4
1M0.1M
27.Thec.m.f.of clfor spontaneouscellreactionhavingtwo
(a)0.415V
(b) 0.374 V
clectrodes as AgIAg+and CuICu2+containing1MCu2+
(c)0.15V
(d)-0.374V
and1MAg+ ions after the passage of 9.65 ampere current
26.In the cel Ag|Ag+IICu$2$+ containing 1 MCu2+ and 1 M
for 1 hour will:
Ag+ ions.9.65 ampere current is passed for 1hour. Select
(a) increase by 0.010 v(b) decrease by 0.010V
the correct statement.
(c)increase by 0.020 V(d)decrease by 0.020 V
(1)Ag will oxidise to Ag+ and new [Ag+] = 1.36 M
(2)Ag+ will reduce to Ag and new [Ag+] = 0.64 M
28.The total change infree energy during the cell reaction in
(3)Cu$2$+ willreduce to Cu and new [Cu$2$+]=0.82 M
CuICu2+IlAg+IAg would be:
(4)Cu will oxidise to Cu2+ and new[Cu$2$+]=0.82 M
(a) -80.1 kJ
(b)-40.05 kJ
(c)-31.36 kJ
(d)-62.72 kJ
AnswersKey
SingleCorrectAnswerType
IntegerType
\end{problembox}

\begin{problembox}
\textbf{Problem 45} \hfill \textbf{[Cengage DPP Q14]}
\label{q:ext-cng_dpp_all_questions-45}

(d)
ComprehensionType
Multiple CorrectAnswersType
\end{problembox}

\begin{problembox}
\textbf{Problem 46} \hfill \textbf{[Cengage DPP Q19]}
\label{q:ext-cng_dpp_all_questions-46}

(a,b, c)
MatchingColumnType
(b$/leftarrow$)(s$/leftarrow$)(1$/leftarrow$P)(A$/leftarrow$)(d$/leftarrow$q)(1$/leftarrow$)
(g$/to$v,u)
--- PAGE 11 ---
Electrochemistry3.11
DPP 3.5
Nernst.Equation and Its Applications
Subjective Type
(C)
(P)
Ecell
1.How will the pH of brine (aq.NaCI solution) be affected
Ecell
when it is electrolysed?
2.From the standard cell potentials given below,calculate
log [Mg2+]
the indicated cell potential E$/circ$.
log
[Zn2+
[Zn2+]
Br(l)
8.The overall cell reaction of Daniell cell is:
E$/circ$
An(s)+ Cu2+Zn2++Cu(s); E$/circ$=+1.10 volt
3.Write the Nernst equationfor the cell reaction in the Daniell
cell. How will the Eeeu be affected when concentration of
[Zn+]
Zn$2$+ ions is increased?
[Cu2+]
(i) How can the reduction potential of an electrode be
increased?
When logio Q is plotted on X-axis and Ecen on Y-axis,
thenwhichofthefollowingwillcorrectlyrepresentthe
Explain why?
variation?
(ii) What is the free energy change for (a) galvanic cell
(b) electrolytic cell?
5.An excess of solid AgCl is added to a 0.1 M solution of
(a)
Brions at 298K.Calculate the concentrations of Cl and
1.10V
Br ions at equilibrium.
Given: E$/circ$(Ci-IAg)= 0.222 V and E(BrlAgBrlAg)=
0.095 V.
6.Find the solubility product of a saturated solution of
AgCrO in water at 298 K if the emf of the cell AglAg+
(satd. Ag2CrO4 soln.)lIAg+ (0.1M))Ag is 0.164 V at
(b)
1.10V
298 K.
SingleCorrectAnswerType
7.In the electrochemical cell,
Mg(s)Mg2(aq)I/Zn$2$t(aq)|Zn(s);E$/circ$= +3.13 V
(c)
1.10V
The correct plot of Ecen versus log !
[Mg$2$+]
will be
represented as:
[Zn$2$+]
(a)Ecell
(b)
Ecell
(d)
log[Mg2]
1.10V
[Zn$2$+]
log [Mg2+]
[Zn2+]
--- PAGE 12 ---
3.12Physical Chemistry
Ag (s) I AgI (saturated solution) Il Ag2C204 (saturated
9.An alloy of Pb-Ag weighing 1.08 g was dissolved in dilute
solution)IAg(s)
HNOand the volume made to100mL.A silver electrode
[Ag+]AgI
was dipped in the solution and the emf of the cell set-up,
(a)Ece= 0.0592 log
Pt(s),H(g)IH*(1M)IAgt(aq)IAg(s) was 0.62 V. If Ecell
[Ag+]Ag2C2O4
is 0.80 V, what is the percentage of Ag in the alloy?
(At 25$/circ$C,RT/F =0.06)
[Ag+]Ag2C2O4
(b) 2.50
(b) Ec==0.0592 log
(a) 25
[Ag+]AgI
(c)10
(d) 50
Ksp(Ag2C204)
10.Determine the solubility and solubility product of
(c) Ecen = 0.0592 log
Co[Fe(CN))] in water at 25$/circ$C from the following data:
Ksp(Ag I)
Conductivity of saturated solution of Co[Fe(CN)]
= 2.06 $/times$ 10-6 $/Omega(-1)$ cm- and that of water = 4.1 $/times$ 10-7
[2Ksg(Ag2C204)]13
$/Omega(-1)$ cm-1. The ionic molar conductivities of Co2+ and
(d)Ecen = 0.0592 log
[Fe(CN)] are 86 and 444 $/Omega$- cm$2$mol-l respectively.
[Ksp(Ag1)]/2
(a)7.69 $/times$ 10-17
(b)2.67$/times$ 10-6
(c)9.68 $/times$10-6
(d) 8.23 $/times$10-17
\end{problembox}

\begin{problembox}
\textbf{Problem 47} \hfill \textbf{[Cengage DPP Q16]}
\label{q:ext-cng_dpp_all_questions-47}

Consider the cell:Ag IAgC1 (s)IKC1 (1M)|Hg2Cl(s)|
\end{problembox}

\begin{problembox}
\textbf{Problem 48} \hfill \textbf{[Cengage DPP Q11]}
\label{q:ext-cng_dpp_all_questions-48}

The standard reduction potential for the following two
Hg (l) I Pt. The cell potential:
reactions are given:
(a)increases on increasing concentration of Cl- ions
AgC1+e$/to$Ag(s)+Cl-(aq);
E$/circ$=0.22V
(i)
(b) decreases or decreasing concentration of Cl- ions
Ag+(ag) +e$/to$Ag(s)
E. =0.80 V
(i)
(c)isindependent of concentration of Cl-ions
The solubility product of AgCl under standard condition
(d) is independent of amounts of AgCl (s) and HgCl (s)
will be:
(a)1.613$/times$10-5M2
(b)1.535$/times$10-8 M2
7.The temperature coefficient of cell is
DE
Then:
(c)3.213$/times$10-10 M2
(d)1.535$/times$10-10M2
aT
\end{problembox}

\begin{problembox}
\textbf{Problem 49} \hfill \textbf{[Cengage DPP Q12]}
\label{q:ext-cng_dpp_all_questions-49}

Calculate the cell potential of following cell:
DE)
When
=0,then $/Delta$H=-nFE
Pt(s)H(g)(0.1 bar)IBOH(0.1 M aq)IHA(0.1 M)IH,(g)
aT
(1 bar)IPt
(b)When
3E
Given : Ka=(HA)= 10-7,K, (BOH)= 10-5
>0,the nFE>$/Delta$H anditis endothermic
aT
(a)0.40V
(b) 0.36V
reaction
(c)0.93V
(d)0.63V
\end{problembox}

\begin{problembox}
\textbf{Problem 50} \hfill \textbf{[Cengage DPP Q13]}
\label{q:ext-cng_dpp_all_questions-50}

For a cell reaction involving two electron charges, the
(c)When
DE
standardemf of the cellis0.295Vat25$/circ$$/circ/text(C)$.Theequilibrium
<0,the nFE<$/Delta$H and it is endothermic
constant of the reaction at 25$/circ$C will be:
dT
reaction
(a)29.5 $/times$10-2
(b)10 
(c)1$/times$110
(d)2.95$/times$10-10
(d)When
dE
=0,then $/Delta$H>nFE and it is endothermic
MultipleCorrectAnswersType
reaction
14.A sample of water has a hardness expressed as 77.5 ppm
18.In the following electrochemical cell:
Ca2+. This sample is passed through an ion exchange
column and the Ca$2$+is replaced by H+.Select the correct
PtlH (x atm)H (pH = y)IZn+ (z M)IZn
statement(s)
Een = Ee.This will be possible when
(a)pH of the water after it has been so treated is 2.4
(b)Every Ca$2$ion is replaced by one H+ ion
(a)[Zn$2$+]=[H]=1 M and pH, = 1 atm
()Every Ca2+ion is replaced by two H+ions
(b) [Zn$2$+]= 0.01,[H]=0.1 M and pH, =1atm
(d)pH of the solution remains unchanged
wne 100="Hd pue NI0=[Hl'W1=I+uZ] (3)
15.Whichof the following expression(s)represent the voltage
ofcell at 298K?
(d)[Zn$2$+]=[H]= 0.1 M and PH = 0.1atm
--- PAGE 13 ---
Electrochemistry3.13
MatchingColumnType
19.Match the Column I with Column II:
ColumnII
(Potential of overall process)
ColumnI
(Combination of half-cell reactions)
(p)
E$/circ$= +1.56V
E=-0.61 V
(a)
6OH-+BrBrO+3HO+6e;
E$/circ$ = -0.76V
7+0H+-O+-HO
E$/circ$=?
++O+O
(b)
E=-0.535V
E=0.45V
(b)
HSO+4H++4eS+3HO;
E$/circ$ =0.17V
SO$2$+4H++2e$/to$HSO+HO;
E$/circ$=?
SO$2$+8H++6eS+4HO; 
E$/circ$=+1.45V
()
E$/circ$=+0.36V
(c)
CIO+6H++6eCI-+3HO;
E$/circ$=+1.36V
Cl+2e2Cl;
CIO+6H++5e-
CI+3H,O;
E$/circ$=?
2
E$/circ$ =-0.76V
(s)
E$/circ$= +1.47 V
(d)
Zn$2$+(aq)+2e---$/to$Zn(s);
Ag++e---$/to$Ag(s);
E=+0.80V
Zn(s)+ 2Ag+Zn$2$+(aq)+2Ag(s);
\end{problembox}

\begin{problembox}
\textbf{Problem 51} \hfill \textbf{[Cengage DPP Q22]}
\label{q:ext-cng_dpp_all_questions-51}

For the following reaction,
IntegerType
\end{problembox}

\begin{problembox}
\textbf{Problem 52} \hfill \textbf{[Cengage DPP Q20]}
\label{q:ext-cng_dpp_all_questions-52}

The voltage of the cell:
Determine the value of In Keq
Pb I PbSO4INa,SO4.10H,O (salt) I Hg,SO41Hg
(a) 58.38
(b) 46.2
(c) 28.3
(d) 66.13
is +0.9647 at 25$/circ$C.The temperature coefficient is 1.74x
23.With-the increase in concentration of NH,pH increases
10V.K-1. Calculate the value of $/Delta$S
to 11.Consider there is no change in concentration of
\end{problembox}

\begin{problembox}
\textbf{Problem 53} \hfill \textbf{[Cengage DPP Q21]}
\label{q:ext-cng_dpp_all_questions-53}

K of Cu(OH) is 1 $/times$ 10-19 If reduction potential of Cu2+/
CH2OandCHO,on addition of NH
Cu couple is 0.1335 V in a solution and E$/circ$ for Cu2+/Cu is
Which of the following is correct?
0.34 V,the pH of solution is
Comprehension Type
(b)Eox decreases by 0.65 V from E
(c) Erea increases by 0.65 Vfrom Eed
Ox
Tollen's reagent is used for the detection of aldehyde when
a solution of AgNO,is added to glucose with NHOH, then
(d) Ered decreases by 0.65 V from Ered
gluconic acid is formed.
24.Ammonia is always added in this reaction. Which of the
Ag++e$/to$Ag
E$/circ$=0.8V
following must be incorrect?
Z+ +HZ+OHO$/circ$H+0H
(a)NH, easily forms complex with Ag+
Eoxi =-0.05 V
(b)Ag(NH)*+ is a strong oxidising reagent than Ag+.
Ag(NH)++e $/to$Ag (s)+2NH
E$/circ$= 0.337V
(c) In absence of NH, silver salt of gluconic acid is.
formed.
(d)NH, has affected the standard reduction potential of
glucose/gluconic acid electrode.
--- PAGE 14 ---
3.14Physical Chemistry
AnswersKey
SubjectiveType
Matching ColumnType
\end{problembox}

\begin{problembox}
\textbf{Problem 54} \hfill \textbf{[Cengage DPP Q2]}
\label{q:ext-cng_dpp_all_questions-54}

0.522V 5.0.0993 M,0.0007 M
\end{problembox}

\begin{problembox}
\textbf{Problem 55} \hfill \textbf{[Cengage DPP Q19]}
\label{q:ext-cng_dpp_all_questions-55}

(a $/to$q),(b $/to$r),(c $/to$s),(d$/to$p)
6.2.29 $/times$ 10-12 M3
Integer Type
SingleCorrectAnswer.Type
\end{problembox}

\begin{problembox}
\textbf{Problem 56} \hfill \textbf{[Cengage DPP Q13]}
\label{q:ext-cng_dpp_all_questions-56}

(c)
ComprehensionType
MultipleCorrectAnswersType
\end{problembox}

\begin{problembox}
\textbf{Problem 57} \hfill \textbf{[Cengage DPP Q24]}
\label{q:ext-cng_dpp_all_questions-57}

(d)
14.(a,c)15. (b,d) 16. (c,d) 17. (a,b,c)
\end{problembox}